\section{Executive summary}

Open Scholar Ranking (OSR) is a free, no-login website that ranks migration destinations for scholars \cite{osr}. It exists because of a decision I had to make myself. When I shortlisted countries for a direct PhD, every ranking told me which university is prestigious, but none told me whether people who moved there stayed. OSR answers that question. It aggregates public ORCID career trajectories \cite{orcid2025}, conditioned on origin nationality, into three metrics at country, city, and institution level: \textbf{Arrival}, a confidence-weighted count of people who reached a destination; \textbf{Stay}, the share who did not later leave that country; and \textbf{No return}, the share who never went back to their home country. A recursive tertile split over the latter two assigns each institution a tier from R1 to R4. I built and deployed the site alone at osr.clarelab.moe.

I benchmarked eight products a prospective mover can consult today, including QS, THE, LinkedIn Alumni, and 1point3acres. None ranks destinations by observed migration conditioned on origin, which is the gap OSR fills. I ran a PACT analysis grounded in 431 archived launch comments, coded into 17 themes, with three interviews with Chinese-speaking students triangulating it. From these I derived 17 functional and 7 non-functional requirements under MoSCoW priorities. Two design--prototype--evaluate iterations follow. The first turned a concept image into the deployed desktop site, evaluated in the field by the public launch and by a wireframe legibility check. The second was a small-screen and touch redesign, tested in a medium-fidelity usability session with three participants, five tasks, and the System Usability Scale \cite{brooke1996sus}. All human data comes from combined sessions, one recorded call per participant: a short interview, the wireframe check, then the usability test.

The launch on rednote and Reddit on 11 June 2026 accumulated roughly 40{,}000 likes and 35{,}000 saves, with 15{,}000 unique site visitors in the surrounding month. The rednote post was taken down and re-uploaded six times in ten days. The comment corpus exposed the launch's central failure. First contact happened on phones, and the site assumed a desktop pointer. Within two weeks the feedback drove a touch-first redesign of every page (Section~\ref{sec:iteration2}). The medium-fidelity test returned a SUS mean of 89.2 with every navigation task completed. It also exposed one weak layer. Two of three participants misread a metric's semantics, which makes point-of-use definitions the next iteration's top fix.

One caveat governs how to read the evidence. This is a solo project. I am designer, developer, and analyst at once. Evidence about the launch is retrospective and externally dated: the git history, the published posts, and the archived comments, with the 58 quotes in the coded themes file verified verbatim against the archives. Evidence produced for this report, the interviews and both evaluation sessions, was collected under the shared protocols. All supporting materials sit in the shared Drive folder, organised by stage and referenced from the relevant sections below.
